Write $N=2^n=[B_n:\Q]$, $q=2^{n+2}$, $\zeta=e^{2\pi i/q}$, $K=\Q(\zeta)$ and $\pi=1-\zeta$; $B_n=\Q(\zeta+\zeta^{-1})\subset K$ with $[K:B_n]=2$.

\emph{Step 1: $2$ is totally ramified in $B_n$ with residue field $\F_2$.} For odd $a$ with $0<a<q$ the quotient $(1-\zeta^a)/(1-\zeta)=1+\zeta+\dots+\zeta^{a-1}$ lies in $\Z[\zeta]$, and if $ab\equiv1\pmod q$ then $(1-\zeta)/(1-\zeta^a)=(1-\zeta^{ab})/(1-\zeta^a)=1+\zeta^a+\dots+\zeta^{a(b-1)}\in\Z[\zeta]$; so $1-\zeta^a=\pi\eta_a$ with $\eta_a\in\Z[\zeta]^\times$. The $q$-th cyclotomic polynomial is $\Phi_q(X)=X^{q/2}+1=\prod_{a\text{ odd}}(X-\zeta^a)$, hence $\prod_{a\text{ odd}}(1-\zeta^a)=\Phi_q(1)=2$, i.e.\ $2=\eta\,\pi^{\varphi(q)}$ with $\eta$ a unit and $\varphi(q)=2^{n+1}=[K:\Q]$. Since the ramification index of a prime above $2$ is at most $[K:\Q]$, this forces $2\mathcal O_K=\mathfrak P^{[K:\Q]}$ with $\mathfrak P=\pi\mathcal O_K$ the unique prime of $K$ above $2$, and its residue degree is $1$. Let $\mathfrak p=\mathfrak P\cap\mathcal O_n$. Every prime of $\mathcal O_n$ above $2$ lies below $\mathfrak P$, so $\mathfrak p$ is the unique such prime, and ramification indices multiply in the tower $\Q\subset B_n\subset K$: $[K:\Q]=e(\mathfrak P/2)=e(\mathfrak P/\mathfrak p)\,e(\mathfrak p/2)\le[K:B_n]\,[B_n:\Q]=[K:\Q]$, whence $e(\mathfrak p/2)=[B_n:\Q]=N$. Then $2\mathcal O_n=\mathfrak p^N$, the residue degree $f(\mathfrak p/2)$ is $1$ (because $e f\le N$), $\mathcal O_n/\mathfrak p=\F_2$, and $2^t\mathcal O_n=\mathfrak p^{Nt}$ for every $t\ge1$.

\emph{Step 2: the unit group of $\mathcal O_n/2^t\mathcal O_n$ is a $2$-group.} Put $k=Nt\ge1$. The ring $\bar R=\mathcal O_n/\mathfrak p^{k}$ is finite of order $\mathrm N(\mathfrak p^{k})=\mathrm N(\mathfrak p)^{k}=2^{k}$ (multiplicativity of the absolute norm in the Dedekind domain $\mathcal O_n$), it is local with unique maximal ideal $\bar{\mathfrak p}=\mathfrak p/\mathfrak p^{k}$ (every ideal of $\bar R$ is $\mathfrak p^i/\mathfrak p^{k}$, $0\le i\le k$), and $|\bar{\mathfrak p}|=\mathrm N(\mathfrak p)^{k-1}=2^{k-1}$. In a finite local ring the units are exactly the elements outside the maximal ideal, so $|\bar R^\times|=2^{k}-2^{k-1}=2^{k-1}$.

\emph{Step 3: the odd power map on $\bar R^\times$ is a bijection.} \lbl{F}: \path{WeberOddTransfer.pow_odd_eq_one_iff_of_two_pow_card} (a group $G$ with $|G|=2^{k'}$ and $\ell$ odd satisfy $g^\ell=1\iff g=1$), \path{pow_odd_eq_iff_of_two_pow_card} ($g^\ell=h^\ell\iff g=h$) and the ring form \path{isUnit_pow_odd_eq_one_iff}; the proof is Mathlib's \path{powCoprime}: when $\gcd(|G|,\ell)=1$ the map $g\mapsto g^\ell$ is a bijection of $G$ with inverse $g\mapsto g^{b}$, $b\ell\equiv1\pmod{|G|}$. Here $|\bar R^\times|=2^{k-1}$ and $\ell$ is odd.

\emph{Step 4: transfer.} Let $x\in\mathcal O_n^\times$ and let $\bar x$ be its image in $\bar R=\mathcal O_n/2^t\mathcal O_n$; $\bar x\in\bar R^\times$ because $x$ is a unit, and for any $y\in\mathcal O_n$ one has $y\equiv1\pmod{2^t\mathcal O_n}$ iff $\bar y=1$. Hence $x^\ell\equiv1\pmod{2^t\mathcal O_n}\iff\bar x^{\ell}=1\iff\bar x=1\iff x\equiv1\pmod{2^t\mathcal O_n}$, the middle equivalence being Step~3. For units $x,y$ the same argument applied to $xy^{-1}\in\mathcal O_n^\times$ gives $x^\ell\equiv y^\ell\iff x\equiv y\pmod{2^t\mathcal O_n}$.

\emph{Remark (valuation form of the same fact).} Since $\mathcal O_n/\mathfrak p=\F_2$, every unit satisfies $x\equiv1\pmod{\mathfrak p}$, so $1+x+\dots+x^{\ell-1}\equiv\ell\equiv1\pmod{\mathfrak p}$ is a $\mathfrak p$-adic unit and $v_{\mathfrak p}(x^\ell-1)=v_{\mathfrak p}(x-1)+v_{\mathfrak p}(1+x+\dots+x^{\ell-1})=v_{\mathfrak p}(x-1)$; the two congruences of the lemma are the two conditions $v_{\mathfrak p}(\cdot)\ge Nt$ on these equal quantities. This second proof is recorded for the reader; the proof above is the one whose discrete core is kernel-checked.
